\section{Độ rộng tích, sai số và chuẩn hóa phép nhân}

Đặt $Y = 4096 + F_Y$ là significand nguyên 13 bit (gồm bit ẩn và 12 bit fraction
cao) của toán hạng được dịch, và $M_0 = Y 2^{12}$ là giá trị nền khởi tạo trên
lưới \texttt{Q24}. Vì $M_0$ có $12$ bit thấp bằng $0$, mọi lượng dịch phải
$S_i \in \{1, \dots, 12\}$ đều bảo toàn toàn bộ trọng số khác $0$ của từng số
hạng tích riêng phần. Khi xử lý trọn vẹn mọi bit $1$ trong 12 bit fraction của $X$,
tích nguyên bằng $(4096+F_X)(4096+F_Y)$ với cận trên $8191^2 = 67\,092\,481 < 2^{26}$.
Do đó 26 bit không dấu của thanh ghi tích lũy \texttt{Q24} (2 bit nguyên và 24 bit
fraction) đủ chứa mọi tích có thể xảy ra mà không tràn ngoài miền hợp lệ; bit 25
bằng $1$ chính là tín hiệu phát hiện tràn chuẩn hóa ($M/2^{24} \ge 2$).

Gọi $a = 1+f_A$ và $b = 1+f_B$ là hai significand đầy đủ ban đầu, còn
$a' = a - \delta_a$ và $b' = b - \delta_b$ là hai significand sau khi giữ 12 bit
fraction cao ($0 \le \delta_a, \delta_b < 2^{-12}$). Độ lệch do bước cắt đầu vào
được phân tách đại số theo biểu thức:
\[
ab - a'b' = a'\delta_b + b'\delta_a + \delta_a\delta_b.
\]
Từ tính chất không âm của các phần dư $\delta_a, \delta_b$ và các số hạng dịch
tiếp theo, với cùng một cặp đầu vào $(A, B)$, ký hiệu độ lớn của tích đầy đủ,
tích 12 bit fraction và tích RTL ở giới hạn $n$ vòng lần lượt là $|P|$, $|P_{12}|$
và $|\widehat P|$, ta luôn có tính đơn điệu theo độ lớn tại đầu ra bộ nhân:
\[
|\widehat P| \le |P_{12}| \le |P|.
\]
Hệ quả quan trọng là bộ nhân luôn mang \textbf{thiên lệch cắt một phía về không}
(\textit{one-sided truncation bias}): phép kiểm tra toàn diện trên toàn bộ
$125.668$ cặp $(A,B)$ hữu hạn khác $0$ của tập mẫu thực nghiệm ở cấu hình
\texttt{OPS1/n12} ghi nhận bất đẳng thức $|\widehat P| \le |P|$ đúng tuyệt đối
trên \textbf{100\% mẫu} ($117.486$ cặp nhỏ hơn chặt và $8.182$ cặp bằng nhau
chính xác khi 14 bit fraction thấp đều bằng $0$). Khi cộng với toán hạng trái dấu
$C$ và đi qua bước lượng tử hóa RNE ở cuối MAC, thiên lệch âm có hệ thống này đôi
khi được bù trừ ngẫu nhiên bởi chiều làm tròn lên của bước đóng gói trung gian
$R(\widehat P)$ trong \texttt{MAC-NF} (giải thích hiện tượng một số phân vị
tương đối đơn bước tại Chương~\ref{ch:measurements}).

\begin{engineeringtable}{Phân tách các nguồn sai số và mốc đối chiếu độc lập.}{tab:error-mechanisms}
\begin{tabularx}{\linewidth}{@{}>{\raggedright\arraybackslash}p{32mm}Y Y@{}}
\toprule
Nguồn sai số & Cơ chế và cờ trạng thái trong RTL & Mốc đối chiếu độc lập \\
\midrule
Cắt fraction đầu vào & Loại bỏ bit thấp ngoài 12 bit ở A hoặc B; chốt cờ \texttt{input\_cut}. & So tích 12 bit $P_{12}$ với tích chính xác $P$ đủ 26 bit fraction. \\
Giới hạn vòng lặp & Mặt nạ $X$ vẫn còn bit 1 khi dừng tại $t=n$; xuất cờ \texttt{approx\_cut}. & So tích RTL $\widehat P$ với $P_{12}$ tại \texttt{core\_done}. \\
Lượng tử hóa khi đóng gói & Giá trị giải mã từ từ mã 32 bit lệch so với giá trị trước làm tròn RNE. & So $\operatorname{decode}(D)$ với $R(\widehat P + C)$ (hoặc $R(R(\widehat P)+C)$). \\
\bottomrule
\end{tabularx}
\end{engineeringtable}

\subsection{Ví dụ minh họa phép nhân Posit(32,3) từng bước}

Xét phép nhân hai số dương $A = 3{,}5 = +2^1 \times 1{,}75$ (mã \texttt{0x47000000})
và $B = 0{,}75 = +2^{-1} \times 1{,}5$ (mã \texttt{0x3E000000}). Trình tự 4 bước
thực hiện phép nhân Posit được minh họa chi tiết như sau:

\begin{enumerate}
  \item \textbf{Bước 1 (Giải mã đầu vào):}
        $A$ cho $s_A = 0$, $sf_A = 1$, $M_A = 1.1100_2$;
        $B$ cho $s_B = 0$, $sf_B = -1$, $M_B = 1.1000_2$.
  \item \textbf{Bước 2 (Tính dấu và scale thô):}
        $s_P = s_A \oplus s_B = 0$ (dương) và $sf_{\mathrm{raw}} = sf_A + sf_B = 1 + (-1) = 0$.
  \item \textbf{Bước 3 (Nhân phần trị theo cột nhị phân):}
        Thực hiện nhân hai significand trên lưới bốn vị trí fraction:
        \[
        \begin{array}{r}
          \phantom{\times}01.1100\\
          \underline{\times\phantom{0}01.1000}\\
          \phantom{+}00.0000\\
          \phantom{+}00.1110\\
          \underline{+\phantom{0}01.1100}\\
          \phantom{+}10.1010
        \end{array}
        \]
        Thu được tích thô $M_{\mathrm{raw}} = 10.1010_2 = 2{,}625_{10}$.
  \item \textbf{Bước 4 (Chuẩn hóa sau nhân và đóng gói):}
        Vì $M_{\mathrm{raw}} = 10.1010_2 \ge 2$ (bit nguyên tràn sang vị trí $2^1$),
        dịch phải 1 bit để đưa bit 1 cao nhất về vị trí ẩn: $M_P = 1.0101_2$ (bỏ bit
        ẩn còn fraction \texttt{0101}), đồng thời tăng scale: $sf_P = sf_{\mathrm{raw}} + 1 = 1$.
        Khi đóng gói, $sf_P = 1 \Rightarrow k = \lfloor 1/8 \rfloor = 0$ (regime
        \texttt{10}) và $e = 1 = \texttt{001}_2$. Ghép trường thu được chuỗi bit
        \texttt{0\_10\_001\_01010...} $= \texttt{0x45400000}$, biểu diễn đúng giá
        trị $2^1 \times 1{,}3125 = 2{,}625_{10}$.
\end{enumerate}

\subsection{Ví dụ khai triển tích lặp trên lưới FULL\_TERM26/Q24}

Để minh họa trực tiếp cơ chế dịch--cộng theo bit trong RTL của đề tài, xét hai
significand $M_X = 1.1010_2 = 1{,}625$ và $M_Y = 1.1000_2 = 1{,}5$. Fraction của
$X$ chỉ có hai bit $1$ tại các vị trí $S_1 = 1$ và $S_2 = 3$ sau dấu chấm nhị phân.
Với quy ước 12 bit fraction, $Y = 4096 + 2048 = 6144$ và giá trị khởi tạo lưới
Q24 là $M_0 = 6144 \times 4096 = 25\,165\,824$. Khai triển tuần tự theo vòng lặp:
\[
1.1010_2 \times 1.1000_2 = 1.1000_2 + 0.1100_2 + 0.0011_2 = 10.0111_2 = 2{,}4375_{10}.
\]
Ở cấp số nguyên Q24 của phần cứng:
\begin{itemize}
  \item Khởi tạo (\texttt{first}, đóng góp bit ẩn): $M_0 = 25\,165\,824$ (tương ứng $M_0/2^{24} = 1{,}5$ nếu $n=0$).
  \item Vòng 1 ($S_1 = 1$): $T_1 = 6144 \times 2^{12-1} = 12\,582\,912 \Rightarrow M_1 = M_0 + T_1 = 37\,748\,736$ (tương ứng $2{,}25$ nếu $n=1$).
  \item Vòng 2 ($S_2 = 3$): $T_2 = 6144 \times 2^{12-3} = 3\,145\,728 \Rightarrow M_2 = M_1 + T_2 = 40\,894\,464$, suy ra $\frac{M_2}{2^{24}} = 2{,}4375$. Mặt nạ của $X$ cạn về $0$ nên lõi tự động dừng sớm ở $t=2$ vòng mà không chạy các vòng rỗng khi cấu hình $n=12$.
\end{itemize}

\paragraph{Đối chiếu trên cặp mã 32 bit FT-05.}
Với bộ dữ liệu kiểm thử RTL \texttt{FT-05}: $A = \texttt{0x1c900000}$ (regime
\texttt{001}, $k_A=-2$, $e_A=6 \Rightarrow sf_A=-10$, $F_A = \texttt{0x480}$) và
$B = \texttt{0x3c820000}$ (regime \texttt{01}, $k_B=-1$, $e_B=7 \Rightarrow sf_B=-1$,
$F_B = \texttt{0x208}$). Cả hai mã đều có 2 bit 1 trong 12 bit fraction cao và
14 bit thấp bằng 0, nên \texttt{OPS1} giữ $A$ làm $X$ và $B$ làm $Y$, cho $b = sf_A + sf_B = -11$.
Quá trình cộng trên lưới Q24 diễn ra như Bảng~\ref{tab:ft05-worked}. Giá trị trước
Packer là $24\,224\,768 \times 2^{-35}$; vì bit 25 bằng 0 nên adapter giữ nguyên
$sf = -11$. Khi $C=0$, cả mô hình tham chiếu và RTL cùng xuất mã \texttt{0x1ae34800}
(khi đặt $n=0$, bộ nhân chỉ phát đóng góp cơ số $M_0 = 18\,907\,136$ và bật cờ
\texttt{approx\_cut=1} để cho ra mã \texttt{0x1a410000}).

\begin{engineeringtable}{Phép tính FULL\_TERM26/Q24 theo hệ số nguyên cho cặp FT-05.}{tab:ft05-worked}
\begin{tabularx}{\linewidth}{@{}>{\raggedright\arraybackslash}p{34mm}r Y@{}}
\toprule
Đại lượng & Giá trị nguyên & Diễn giải nguồn gốc \\
\midrule
$Y$ & 4.616 & Bit ẩn $4.096$ cộng fraction $B$ ($F_B = 520$). \\
$M_0$ & 18.907.136 & Cơ số $Y \ll 12$ tại xung \texttt{first}. \\
Term 1 ($S_1=2$) & 4.726.784 & Bit 1 đầu tiên của $X$, $M_0 \gg 2$. \\
Term 2 ($S_2=5$) & 590.848 & Bit 1 thứ hai của $X$, $M_0 \gg 5$. \\
$M$ tại \texttt{last} & 24.224.768 & Tổng cơ số và hai term, $t=2$, $b=-11$. \\
\bottomrule
\end{tabularx}
\end{engineeringtable}

\section{Nguyên lý cộng, căn chỉnh, jam và chuẩn hóa triệt tiêu}

Khi cộng hoặc trừ hai số hạng có dấu, scale và significand đã chuẩn hóa, bước
đầu tiên là so sánh độ lớn (so $sf$ trước, nếu bằng nhau thì so significand) để
đặt số hạng lớn hơn làm toán hạng nền $X$ và số hạng còn lại làm $Y$. Khi hai
toán hạng cùng dấu hiệu dụng, thực hiện phép cộng $M_X + (M_Y \gg \Delta sf)$;
khi khác dấu hiệu dụng, thực hiện phép trừ không âm $M_X - (M_Y \gg \Delta sf)$
và giữ dấu của $X$.

\paragraph{Ví dụ minh họa phép cộng phát sinh carry và phép trừ triệt tiêu.}
Xét cặp toán hạng $A = 3{,}5$ (\texttt{0x47000000}, $sf_A = 1$, $M_A = 1.1100_2$)
và $B = 0{,}75$ (\texttt{0x3E000000}, $sf_B = -1$, $M_B = 1.1000_2$):
\begin{enumerate}
  \item \textbf{Phép cộng có nhớ ($A + B$):} Vì $sf_A > sf_B$, chọn $A$ làm nền
        với độ lệch $\Delta sf = 1 - (-1) = 2$. Dịch phải $M_B$ hai vị trí thu
        được $0.0110_2$ trên cùng thang $2^1$:
        \[
        3{,}5 + 0{,}75 = 2^1\bigl(1.1100_2 + 0.0110_2\bigr),
        \qquad
        \begin{array}{r}
          \phantom{+}01.1100\\
          \underline{+\phantom{0}00.0110}\\
          \phantom{+}10.0010
        \end{array}
        \]
        Bit nhớ tràn vào vị trí nguyên thứ hai tạo tổng thô $10.0010_2$. Bước chuẩn
        hóa sau carry dịch phải 1 bit và tăng scale thêm $1$:
        $10.0010_2 \times 2^1 = 1.0001_2 \times 2^2 = 4{,}25$. Mã hóa $k=0$, $e=2$
        (\texttt{010}), fraction \texttt{0001} cho ra mã kết quả \texttt{0x48400000}.
  \item \textbf{Phép trừ gây triệt tiêu ($A - 3{,}25$):} Thay số hạng thứ hai
        bằng $-3{,}25$ ($s=1$, $sf=1$, $M = 1.1010_2$). Do cùng $sf=1$ và khác
        dấu, phép cộng chuyển thành trừ hai độ lớn:
        \[
        \begin{array}{r}
          \phantom{-}01.1100\\
          \underline{-\phantom{0}01.1010}\\
          \phantom{-}00.0010
        \end{array}
        \qquad\Longrightarrow\qquad
        0.0010_2 \times 2^1 = 1.0000_2 \times 2^{-2} = 0{,}25.
        \]
        Hiệu độ lớn mất ba bit cao, buộc mạch chuẩn hóa phải dịch trái $3$ vị
        trí và giảm scale từ $1$ xuống $-2$, đóng gói thành \texttt{0x38000000}.
        Trong cửa sổ cộng 31 bit có vị trí bit ẩn $h=29$, ký hiệu $j$ là chỉ số
        bit 1 cao nhất của tổng/hiệu và $\ell = 30 - j$ là số bit 0 dẫn đầu đếm
        từ bit 30, công thức cập nhật scale hợp nhất cho cả trường hợp có nhớ
        ($\ell = 0$), bình thường ($\ell = 1$) và triệt tiêu ($\ell > 1$) là:
        \[
        sf_{\mathrm{out}} = sf_{\mathrm{base}} + j - h = sf_{\mathrm{base}} + 1 - \ell.
        \]
        Áp dụng cho ví dụ trừ trên với $sf_{\mathrm{base}}=1$ và $\ell=4$ ($j=26$)
        thu được đúng $sf_{\mathrm{out}} = 1 + 1 - 4 = -2$. Nếu hai độ lớn triệt
        tiêu hoàn toàn bằng $0$, cờ \texttt{Zero} được kích hoạt trực tiếp thay
        cho kết quả đếm của LOD.
\end{enumerate}

\paragraph{Phần dư căn chỉnh và thao tác jam.}
Trong cửa sổ cộng độ rộng hữu hạn, các bit của toán hạng nhỏ $Y$ bị dịch ra ngoài
cạnh phải của cửa sổ được gom bằng phép OR (\textit{sticky}). Nếu giữ bit LSB
của cửa sổ phản ánh phần dư khác 0 (gọi là \textit{jam}), mạch biết được độ lớn
thật của số hạng nhỏ lớn hơn phần còn giữ trong cửa sổ. Trong phép trừ, phần dư
này làm giảm nhẹ độ lớn hiệu và ảnh hưởng trực tiếp tới biên làm tròn khi xảy ra
triệt tiêu; vì vậy đường dữ liệu phải truyền thông tin phần dư qua cả tầng số học
và tầng chuẩn hóa đến Packer thay vì chỉ OR mù một cờ mất dữ liệu ở cuối.

\section{Làm tròn RNE, mã hóa đầu ra và ranh giới ba kiến trúc}

Sau chuẩn hóa, từ $sf_{\mathrm{out}}$ tách $k = \lfloor sf_{\mathrm{out}} / 8 \rfloor$
và $e = sf_{\mathrm{out}} - 8k \in [0, 7]$ bằng phép chia lấy phần nguyên xuống
(ví dụ $sf = -11 \Rightarrow k = -2, e = 5$). Chuỗi bit $1+\text{terminator}$ hoặc
$0+\text{terminator}$ của $k$ ghép với exponent và fraction tạo thành vector độ
lớn được giữ lại $31$ bit. Gọi $L$ là bit thấp nhất còn giữ, $G$ (\textit{Guard})
là bit kề ngay sau $L$, $R$ (\textit{Round}) là bit tiếp theo và $S$ (\textit{Sticky})
là phép OR toàn bộ các bit thấp hơn cùng phần dư căn chỉnh, điều kiện tăng mã
độ lớn theo quy tắc RNE (\textit{Round to Nearest, Ties to Even}) là:
\[
\rho = G \land (L \lor R \lor S).
\]

\begin{engineeringtable}{Bảng quyết định làm tròn RNE theo các bit L, G, R và S.}{tab:rne-decision}
\begin{tabularx}{\linewidth}{@{}llllY@{}}
\toprule
$G$ & $R$ & $S$ & $L$ & Hành động làm tròn \\
\midrule
0 & bất kỳ & bất kỳ & bất kỳ & Làm tròn xuống (giữ nguyên độ lớn). \\
1 & 0 & 0 & 0 & Hòa chẵn (\textit{tie-even}): giữ mã có $L=0$. \\
1 & 0 & 0 & 1 & Hòa lẻ (\textit{tie-odd}): tăng 1 đơn vị để chuyển $L$ về $0$. \\
1 & 1 & bất kỳ & bất kỳ & Vượt nửa bước: tăng độ lớn thêm 1 đơn vị. \\
1 & 0 & 1 & bất kỳ & Vượt nửa bước (do có dư ở $S$): tăng độ lớn thêm 1 đơn vị. \\
\bottomrule
\end{tabularx}
\end{engineeringtable}

\paragraph{Ví dụ hợp nhất làm tròn RNE và mã hóa bù hai cho số âm.}
Giả sử phần độ lớn giữ được đang là $m = 6$ và quyết định RNE là tăng $\rho = 1$.
Nếu kết quả dương, đầu ra là $m + \rho = 7$; nếu kết quả mang dấu âm, giá trị cần
xuất là $-(m+\rho) = -7$. Trên từ mã 32 bit, đồng nhất thức đại số
\[
-(m + \rho) = \overline{m} + (1 - \rho) \pmod{2^{32}}
\]
cho thấy thay vì cộng xong rồi mới lấy bù hai bằng hai bộ cộng nối tiếp, phần
cứng chỉ cần đảo bit $m$ theo dấu kết quả và đưa bit cộng đầu vào bằng
$s \oplus \rho$: cụ thể $\overline{6} + (1 - 1) = \overline{6} = -7$ (khi $\rho=1$),
còn khi $\rho=0$ thì $\overline{6} + (1 - 0) = -6$. Các trường hợp ngoại lệ
\texttt{Zero}, \texttt{NaR} hoặc kẹp biên $\pm\mathrm{minpos}/\pm\mathrm{maxpos}$
được xử lý song song qua bộ chọn ở ngõ ra.

\paragraph{Hệ quả của vị trí làm tròn đối với MAC-NF, MAC-F và MAC-Quire.}
Đặt $\varepsilon_p = R(\widehat P) - \widehat P$ là sai số làm tròn khi đóng gói
tích về \texttt{Posit(32,3)}, quan hệ giữa giá trị đầu vào bộ cộng của hai kiến
trúc là $R(\widehat P) + C = \widehat P + C + \varepsilon_p$. Khi $\widehat P$ và
$C$ trái dấu và gần triệt tiêu nhau, tổng thực rất nhỏ nên $\varepsilon_p$ trở thành
nhiễu đáng kể làm sai lệch chữ số có nghĩa hoặc ngưỡng kẹp (như đã đo được trên nhóm
$C=R(-P)$ ở Chương~\ref{ch:measurements}). Kiến trúc \texttt{MAC-F} loại bỏ được
$\varepsilon_p$ ở từng phép MAC; tuy nhiên trong chuỗi tích lũy nhiều bước hồi tiếp
ngoài, trạng thái $C$ vẫn phải đi qua một lần lượng tử hóa ở cuối mỗi giao dịch,
trong khi \texttt{MAC-Quire} loại bỏ hoàn toàn bước lượng tử hóa trạng thái trung
gian nhờ thanh ghi $Q$ trên lưới $2^{-480}$.
